\section{Classical inputs}\label{sec:classical-inputs}
\raggedbottom

We use counting-measure $\ell^2$ norms. Write $n\sim Y$ for $Y<n\le2Y$.
The symbols $\mu_{\rm Mob},\Lambda,\varphi$ have their usual meanings.
Write $d(n)$ for the divisor function.
The parameters in each fact are local.
Put $P^-(1)=\infty$ and $e(t)=\exp(2\pi it)$.
Statements (P1.x) and (P2.x) restate results proved in
Papers~I and~II, respectively, with their original scope.
The structured sieve result is Paper~I, Theorem~2.6.
Its classical source is Iwaniec's Theorem~4
\cite[Theorem~4]{Iwaniec1980Remainder}. The \(S\)-form displayed below
specializes the \(W_R^\pm\) estimates of Paper~I, Theorem~2.6.
It follows from them by \(S\le W_R^+\) and \(S\ge W_R^-\).
Theorem~4 does not state it verbatim.
Each new application needs its own coefficient, support and
uniformity checks.

\begin{fact}[(F1) Prime estimates]\label{fact:prime-estimates}
For every $A>0$, one may take $B_\psi(A)=A+60$ in
\begin{equation}\label{eq:classical-bv}
 \begin{gathered}
 \sum_{q\le Q}\max_{(a,q)=1}
 \left|\sum_{\substack{n\le y\\n\equiv a\ (q)}}\Lambda(n)
                -\frac y{\varphi(q)}\right|
 \ll_A y(\log y)^{-A},\\
 y\ge3,\quad Q\le y^{1/2}(\log y)^{-B_\psi(A)}.
 \end{gathered}
\end{equation}
For fixed $A,C>0$, uniformly for $q\le(\log y)^C$,
\begin{equation}\label{eq:classical-sw}
 \pi(y;q,a)=\frac{\operatorname{Li}(y)}{\varphi(q)}
       +O_{A,C}(y(\log y)^{-A}),\quad (a,q)=1,\quad y\ge3.
\end{equation}
These are Paper~I, (2.20)--(2.21) \cite[(2.20)--(2.21), p.~9]{PaperI}.
The endpoint $y$ is fixed. The second constant is ineffective.
\end{fact}

\begin{fact}[(F1$'$) Maximal-endpoint prime estimates]\label{fact:prime-estimates-maximal}
Use the prime Siegel--Walfisz estimate \eqref{eq:classical-sw}.
For every fixed $S>0$ and $X\ge3$, put
$Q_X=X^{1/2}(\log X)^{-(S+60)}$. Then
\begin{equation}\label{eq:classical-bv-maximal}
 \begin{aligned}
 \sum_{q\le Q_X}\max_{(a,q)=1}\sup_{0\le y\le X}
 \left|\sum_{\substack{n\le y\\n\equiv a\ (q)}}\Lambda(n)
                -\frac y{\varphi(q)}\right|
 &\ll_S X(\log X)^{-S},\\
 \sum_{q\le Q_X}\max_{(a,q)=1}\sup_{0\le y\le X}
 \left|\pi(y;q,a)-\frac{\operatorname{Li}(y)}{\varphi(q)}\right|
 &\ll_S X(\log X)^{-S}.
 \end{aligned}
\end{equation}
Here $\operatorname{Li}(y)=0$ for $0\le y<2$.
The level is imposed at $X$; the endpoints may be chosen separately for
each modulus. An interval difference costs at most twice the bound.
This is the classical maximal form of Bombieri--Vinogradov
\cite[Chapter~28]{Davenport1980}. Removing prime powers and partial
summation pass from the von Mangoldt form to the prime-counting form.
The constants may be ineffective through \eqref{eq:classical-sw}.
\end{fact}

\begin{fact}[(F2) Full-support bilinear distribution]\label{fact:bv-full-support}
Let $a_m,b_n$ have supports $m\sim M,n\sim Y$.
Assume $x=MY\ge3$, set $L_x=\log x$, and fix $\epsilon,B>0$.
Assume $\epsilon\le\log Y/\log x\le1-\epsilon$.
For every $A_0>0$, assume uniformly in $h,k\ge1$ and integers
$\ell\ne0$ with $(\ell,k)=1$ that
\begin{equation}\label{eq:bvfs-sw}
 \left|\sum_{\substack{n\equiv\ell\ (k)\\(n,h)=1}}b_n
 -\frac1{\varphi(k)}\sum_{(n,hk)=1}b_n\right|
 \le C_{A_0}\|b\|_2Y^{1/2}d(h)^B(\log(2Y))^{-A_0}.
\end{equation}
Define
\[
 \Delta_{\rm fs}(q,a)=\sum_{mn\equiv a\ (q)}a_mb_n
       -\frac1{\varphi(q)}\sum_{(mn,q)=1}a_mb_n.
\]
For every $A>0$, some $B_{\rm BV}=B_{\rm BV}(A,B,\epsilon)>0$ gives
\begin{equation}\label{eq:bvfs}
 \sum_{q\le Q}\max_{(a,q)=1}|\Delta_{\rm fs}(q,a)|
 \ll \|a\|_2\|b\|_2x^{1/2}L_x^{-A},\qquad
 Q\le x^{1/2}L_x^{-B_{\rm BV}}.
\end{equation}
Both sums use the full supports. The constants depend on the family constants
$C_{A_0}$. This follows from the proof of BFI Theorem~0(b)
\cite[(A1)--(A2), p.~206; Theorem~0(b), (1.5), p.~212;
proof, pp.~213--214]{BFI86}.
\end{fact}

\begin{proof}
Character orthogonality and primitive induction give
\[
\begin{aligned}
 \text{left side of \eqref{eq:bvfs}}
 &\le\sum_{e\le Q}\frac1{\varphi(e)}
       \sum_{2\le f\le Q/e}\frac1{\varphi(f)}\\
 &\qquad\cdot\sum_{\chi\ (f)}^*|A_e(\chi)B_e(\chi)|,\\
 A_e(\chi)&=\sum_{(m,e)=1}a_m\chi(m).
\end{aligned}
\]
Define $B_e$ similarly. Split at $F=L_x^K$.
Since $M\ge x^\epsilon$, large $x$ gives $F^2\le M$.
This permits the small-conductor estimate. The hypothesis handles $f\le F$.
The multiplicative large sieve handles $f>F$. Before summing $e$, the bounds are
\[
 \ll\|a\|_2\|b\|_2x^{1/2}d(e)^BF^{3/2}L_x^{-A_0},
 \qquad
 \ll L_x^2\|a\|_2\|b\|_2
       (Q+M^{1/2}+Y^{1/2}+x^{1/2}/F).
\]
Choose $B_2$ with $\sum_{e\le Q}d(e)^B/\varphi(e)\ll L_x^{B_2}$.
In the large-conductor bound, summing $1/\varphi(e)$ costs
one more logarithm, since $\sum_{e\le Q}1/\varphi(e)\ll L_x$.
Take $K=A+B_2+4$, $A_0=A+B_2+3K/2+1$, and $B_{\rm BV}=K$.
Since $M,Y\ge x^\epsilon$, these bounds prove \eqref{eq:bvfs}.
The large sieve is BFI (1.6), p.~207, restated in Paper~I, (2.19)
\cite[(1.6), p.~207]{BFI86}\cite[(2.19), p.~9]{PaperI}.
\end{proof}

We do not use BFI86, Theorem~0(a).

Fact (F2) is the standalone full-support bound.
Section~\ref{sec:l3} checks the coefficients, relative levels and rough
factors of the first three switched rows on the actual source.

\begin{fact}[(F3) The spectral gap]\label{fact:kim-sarnak}
For every congruence subgroup $\Gamma\subset\mathrm{SL}_2(\mathbb Z)$,
\begin{equation}\label{eq:classical-spectral-constant}
 \lambda_1(\Gamma)\ge\frac{975}{4096},\qquad
 0<\lambda=\frac14-t^2<\frac14,\quad t\ge0\Longrightarrow
 0\le t\le\frac7{64},\quad \vartheta=2t\le\frac7{32}.
\end{equation}
Here $\lambda_1(\Gamma)$ is the smallest nonzero Laplace eigenvalue.
The implication applies to every eigenvalue $\lambda$ of $\Delta$
on $L^2(\Gamma\backslash\mathbb H)$ in the indicated range.
The first bound is Kim--Sarnak \cite[Appendix~2]{KimSarnak2003}.
The implications are elementary.
\end{fact}

\begin{fact}[(F4) Deligne's bound]\label{fact:F4-deligne}\label{fact:F5-deligne}
Let $f=\sum_{n\ge1}\psi_f(n)e(nz)$ be a primitive holomorphic cuspidal newform.
Assume $\psi_f(1)=1$, integer weight $k\ge2$, level $M$, and character $\chi$,
with $f((az+b)/(cz+d))=\chi(d)(cz+d)^kf(z)$ on $\Gamma_0(M)$.
This agrees with Deligne's convention $\chi(a)^{-1}$.
For $p\nmid M$, Deligne proved
\begin{equation}\label{eq:deligne-roots}
 |\psi_f(p)|\le2p^{(k-1)/2}
\end{equation}
\cite[Th\'eor\`eme~(8.2), p.~302]{Deligne1974}.
The same source states that the roots $\alpha_p,\beta_p$ of
$T^2-\psi_f(p)T+\chi(p)p^{k-1}$ both have modulus $p^{(k-1)/2}$.
Since $|\alpha_p|^2=|\beta_p|^2=p^{k-1}$, we have
$\overline{\alpha_p}=p^{k-1}/\alpha_p$ and likewise for $\beta_p$.
With $\alpha_p\beta_p=\chi(p)p^{k-1}$, this gives
\begin{equation}\label{eq:deligne-conjugation}
 \begin{aligned}
 \chi(p)\overline{\psi_f(p)}
 &=\frac{\alpha_p\beta_p}{p^{k-1}}
   \left(\frac{p^{k-1}}{\alpha_p}+\frac{p^{k-1}}{\beta_p}\right)
 =\psi_f(p).
 \end{aligned}
\end{equation}
Choose $\eta_p^2=\chi(p)$, $|\eta_p|=1$.
Then $b_p=\eta_p^{-1}p^{-(k-1)/2}\psi_f(p)$ is real and $|b_p|\le2$.
With $T=p^{(k-1)/2}\eta_p Z$, the polynomial becomes
\[
 T^2-\psi_f(p)T+\chi(p)p^{k-1}
       =p^{k-1}\chi(p)(Z^2-b_pZ+1).
\]
Write $b_p=2\cos\omega_p$, $0\le\omega_p\le\pi$.
The normalized roots are $e^{\pm i\omega_p}$,
also in the repeated case $b_p=\pm2$.
An earlier exposition of the modular-form setting is
\cite{DeligneBourbaki1971}.
\end{fact}

\begin{fact}[(F5) Averaged Kloosterman sums]\label{fact:F5-di82-thm14}\label{fact:F2-di82-thm14}
For $C,M,Y\ge1$ and $\epsilon>0$,
\begin{equation}\label{eq:F2-di82-thm14}
 \sum_{c\le C}\left|\sum_{m\le M}\sum_{n\le Y}S(m,n;c)\right|
 \ll_\epsilon(CMY)^\epsilon C(C+MY)
\end{equation}
\cite[Theorem~14, equation~(8.13), p.~275]{DI82}.
The sums run over integers $1\le c\le C$, $1\le m\le M$,
and $1\le n\le Y$. Here $S(m,n;c)=\sum_{u\ (c)}^*e((mu+n\bar u)/c)$.
\end{fact}

Use the trivial character on $\Gamma_0(q)\backslash\mathbb H$.
The sum over $f$ runs through a Petersson-orthonormal basis of
$S_k(\Gamma_0(q))$. The sum over $j$ runs through an orthonormal basis of
Maass cusp forms, including exceptional forms. The sum over $\mathfrak a$
runs through inequivalent cusps with width-one scaling matrices.
The measure is $dx\,dy/y^2$.
At infinity, the DI coefficients satisfy
\begin{equation}\label{eq:di-normalization}
 \begin{aligned}
 u_j(z)&=\sum_{n\ne0}\rho_j(n)\sqrt y
            K_{i\kappa_j}(2\pi|n|y)e(nx),\quad
 \lambda_j=\tfrac14+\kappa_j^2,\\
 f(z)&=\sum_{n\ge1}\psi_f(n)e(nz),\\
 E_{\mathfrak a}(z,s)&=\text{constant terms}
 +\sqrt y\sum_{n\ne0}\frac{2\pi^s|n|^{s-1/2}}{\Gamma(s)}
       \varphi_{\mathfrak a\infty n}(s)
       K_{s-1/2}(2\pi|n|y)e(nx).
 \end{aligned}
\end{equation}
For $0<\lambda_j<1/4$, set
$\vartheta_j=2|\Im\kappa_j|=2\sqrt{1/4-\lambda_j}$.
See Paper~I, (2.2)--(2.4) \cite[(2.2)--(2.4), p.~5]{PaperI}.

\begin{fact}[(F6) Spectral large sieves]\label{fact:di-large-sieves}
For $q,K\ge1$, $Y\ge1/2$, and $a_n$ supported on $n\sim Y$,
put
\[
 \begin{aligned}
 \Sigma_{\rm hol}(K,a)&=\sum_{\substack{2\le k\le K\\k\ {\rm even}}}
 \frac{(k-1)!}{(4\pi)^{k-1}}\sum_f
 \left|\sum_{n\sim Y}a_nn^{-(k-1)/2}\psi_f(n)\right|^2,\\
 \Sigma_{\rm Maa}(K,a)&=\sum_{|\kappa_j|\le K}\frac1{\cosh(\pi\kappa_j)}
            \left|\sum_{n\sim Y}a_n\rho_j(n)\right|^2,\\
 \Sigma_{\rm Eis}(K,a)&=\sum_{\mathfrak a}\int_{-K}^K
 \left|\sum_{n\sim Y}a_nn^{it}
       \varphi_{\mathfrak a\infty n}(1/2+it)\right|^2dt.
 \end{aligned}
\]
The large-sieve estimates are
\begin{equation}\label{eq:di-large-sieves}
 \begin{gathered}
 \Sigma_{\rm hol}(K,a),\Sigma_{\rm Maa}(K,a),\Sigma_{\rm Eis}(K,a)
       \ll_\epsilon(K^2+Y^{1+\epsilon}/q)\|a\|_2^2,\\
 \sum_{j\ {\rm exceptional}}V^{\vartheta_j}
       \left|\sum_{n\sim Y}a_n\rho_j(n)\right|^2
 \ll_\epsilon(1+\sqrt{YV/q})(1+\sqrt{Y^{1+\epsilon}/q})\|a\|_2^2,
 \quad V\ge1.
 \end{gathered}
\end{equation}
These are the specialized forms in Paper~I, (2.5)--(2.8)
\cite[Theorem~2, (1.28)--(1.30), p.~230;
Theorem~5, (1.38), p.~232]{DI82}
\cite[(2.5)--(2.8), pp.~6--7]{PaperI}.
\end{fact}

For $\phi\in C_c^\infty((0,\infty))$, define
\begin{equation}\label{eq:di-transforms}
 \begin{aligned}
 \widehat\phi(t)&=\frac\pi{\sinh\pi t}\int_0^\infty
 \frac{J_{2it}(x)-J_{-2it}(x)}{2i}\phi(x)\frac{dx}x,\\
 \check\phi(t)&=\frac4\pi\cosh\pi t\int_0^\infty
 K_{2it}(x)\phi(x)\frac{dx}x,\qquad
 \widetilde\phi(l)=\int_0^\infty J_l(x)\phi(x)\frac{dx}x.
 \end{aligned}
\end{equation}
The value at $t=0$ is the removable limit.

\begin{fact}[(F7) Corrected same-cusp trace formula]\label{fact:p1-corrected}
For $q,m,n\ge1$ and the test class above,
\begin{equation}\label{eq:p1-corrected}
 \sum_{h\ge1}\frac{S(m,n;qh)}{qh}
 \phi\left(\frac{4\pi\sqrt{mn}}{qh}\right)
       =\frac{H_\phi(m,n)}\pi-M_\phi(m,n)-\frac{E_\phi(m,n)}\pi,
\end{equation}
where
\begin{align*}
 H_\phi(m,n)&=\sum_{\substack{k\ge2\\k\ {\rm even}}}
 \frac{i^k(k-1)!\widetilde\phi(k-1)}{(4\pi\sqrt{mn})^{k-1}}
       \sum_f\overline{\psi_f(m)}\psi_f(n),\\
 M_\phi(m,n)&=\sum_j\frac{\widehat\phi(\kappa_j)}{\cosh\pi\kappa_j}
                  \overline{\rho_j(m)}\rho_j(n),\\
 E_\phi(m,n)&=\sum_{\mathfrak a}\int_{\mathbb R}(n/m)^{it}
 \overline{\varphi_{\mathfrak a\infty m}(1/2+it)}
 \varphi_{\mathfrak a\infty n}(1/2+it)\widehat\phi(t)\,dt.
\end{align*}
No extra diagonal term occurs, also for $m=n$.
This is the corrected same-cusp specialization proved in Paper~I,
(2.7.T15)--(2.7.T18) \cite[(2.7.T15)--(2.7.T18)]{PaperI}. Compare
\cite[Lemma~4.5, (4.8)--(4.17)]{Drappeau2017}, in the DI convention
\cite[(1.15)--(1.17)]{DI82}.
\end{fact}

\begin{fact}[(F8) Mixed-cusp formulas]\label{fact:mixed-cusps}
For $q=rs$, $(r,s)=1$, use the trivial character and
\[
 \sigma_{1/s}=\begin{pmatrix}\sqrt r&0\\s\sqrt r&1/\sqrt r\end{pmatrix},
 \quad\mathcal C(\infty,1/s)=\{cs\sqrt r:c\ge1,\ (c,r)=1\}.
\]
DI82 (1.6), p.~224, gives the arithmetic identity
\cite[(1.6), p.~224]{DI82}
\[
 S_{\infty,1/s}(m,n;cs\sqrt r)
       =e(n\bar s/r)S(m\bar r,n;cs),\qquad m,n\in\mathbb N.
\]
Here $s\bar s\equiv1\pmod r$ and $r\bar r\equiv1\pmod{cs}$.
The accompanying trace formulas, proved in Paper~I, are
\[
 \mathcal K_+=H_{\infty,1/s}/\pi-M_{\infty,1/s}-E_{\infty,1/s}/\pi,
 \qquad
 \mathcal K_-=(\pi/2)M^-_{\infty,1/s}+(1/2)E^-_{\infty,1/s}.
\]
Here $\mathcal K_\pm=\sum_{\gamma\in\mathcal C}\gamma^{-1}
S_{\infty,1/s}(m,\pm n;\gamma)\phi(4\pi\sqrt{mn}/\gamma)$.
The plus terms have the definitions in (F7), with their second
coefficient at $1/s$.
For $M^-$ and $E^-$, replace $\widehat\phi$ by $\check\phi$ and $n$ by $-n$
only in the second Fourier coefficient.
The cusps are singular, and the scaling matrices have width one.
The classical width of $1/s$ is $r$.
For $r>1$ they are inequivalent.
The normalization is that of Paper~I, (2.13)--(2.14). Compare
\cite[Lemmas~4.3, 4.5, (4.7)--(4.19)]{Drappeau2017}
\cite[(2.13)--(2.14), p.~8]{PaperI}.
\end{fact}

\begin{fact}[(F9) Bessel representations]\label{fact:F9-dlmf}\label{fact:F7-dlmf}
For $x>0$, $n\in\mathbb Z$, and $|\Re\nu|<1$,
\begin{align*}
 J_n(x)&=\pi^{-1}\int_0^\pi\cos(x\sin u-nu)\,du\\
       &\quad=i^{-n}\pi^{-1}\int_0^\pi e^{ix\cos u}\cos(nu)\,du,\\
 J_\nu(x)&=(2/\pi)\int_0^\infty
       \sin(x\cosh u-\pi\nu/2)\cosh(\nu u)\,du,\\
 Y_\nu(x)&=-(2/\pi)\int_0^\infty
       \cos(x\cosh u-\pi\nu/2)\cosh(\nu u)\,du,\\
 K_\nu(x)&=\sec(\pi\nu/2)\int_0^\infty
       \cos(x\sinh u)\cosh(\nu u)\,du.
\end{align*}
The infinite integrals are improper.
These are DLMF (10.9.2), (10.9.8), and (10.32.7)
\cite[(10.9.2), (10.9.8), (10.32.7)]{DLMF}.
Paper~I uses them on printed p.~20, (2.7.T31)--(2.7.T32).
\end{fact}

\begin{fact}[(F10) Kontorovich--Lebedev inversion]\label{fact:kl-inversion}
For $g$ on $(0,\infty)$, the inversion pair is
\[
 f(y)=\frac{2y\sinh\pi y}{\pi^2}\int_0^\infty g(x)K_{iy}(x)\frac{dx}x,
 \qquad g(x)=\int_0^\infty f(y)K_{iy}(x)\,dy.
\]
It holds under either of two conditions.
(a) $g/x\in C^1$, with $xg$ and $x(g/x)'$ absolutely integrable.
(b) Local piecewise continuity and bounded variation, with
\[
 \int_0^{1/2}\frac{|g(x)|}x\log(1/x)\,dx<\infty,\qquad
 \int_{1/2}^\infty\frac{|g(x)|}{\sqrt x}\,dx<\infty.
\]
Under condition (b), the inverse equality is asserted at continuity points of $g$.
In our applications, the inverse equality is used only at continuity points
of $g$.
See DLMF (10.43.30)--(10.43.32)
\cite[(10.43.30)--(10.43.32)]{DLMF}.
These hypotheses hold for $g=\phi/x$, $\phi\in C_c^\infty((0,\infty))$.
They also hold for $g=xK_{2it}(x)$ when $|\Im t|\le7/64$.
Near zero, $K_{2it}(x)\ll_t x^{-2|\Im t|}(1+|\log x|)$,
so the first integral in condition (b) converges absolutely.
At infinity, $K_{2it}(x)\ll_t x^{-1/2}e^{-x}$, so the second converges too.
Both functions are smooth on compact subintervals.
\end{fact}

\begin{fact}[(F11) Hecke operators on the full space]\label{fact:hecke-full-space}
For trivial character and $(k,q)=1$, Paper~I proves
\[
 \|T_k\|_{\mathrm{Maass},q}\le d(k)k^{7/64},\qquad
 \|T_k\|_{\mathrm{hol},q}\le d(k).
\]
Their proof uses the orthogonal cuspidal representation
decomposition of Jacquet--Langlands, Chapter~III, \S16
\cite[Chapter~III, \S16]{JacquetLanglands1970}.
It also uses the holomorphic oldform classes of Atkin--Lehner
\cite[Lemmas~6, 13--15; Theorems~3(i), 5]{AtkinLehner1970}.
Atkin--Lehner's weight $k$ means classical weight $2k$.
For an unramified normalized newform, the algebra needed is
\[
 \lambda_f(p^r)=\sum_{j=0}^r\alpha_p^{r-j}\beta_p^j,\qquad
 \lambda_f(mn)=\lambda_f(m)\lambda_f(n)\quad((m,n)=1).
\]
The Maass Satake bound is $|\alpha_p|,|\beta_p|\le p^{7/64}$
\cite[Appendix~2]{KimSarnak2003}.
The holomorphic bound is (F4). Full-space diagonalization and unitary
transport give the displayed norms, as in Paper~I, (2.7.T37)--(2.7.T40)
\cite[(2.7.T37)--(2.7.T40), pp.~21--22]{PaperI}.
\end{fact}

\begin{fact}[(F12) Eisenstein inputs]\label{fact:eisenstein-inputs}
Paper~I uses three classical Eisenstein properties.
The Eisenstein family has a meromorphic continuation, which allows compact
Fourier-coefficient extraction. It has an arithmetic Fourier normalization.
There is a complete irredundant cusp set.
For the cusp $u/w$, let $d_{\rm den}$ be the lower-row denominator variable
in $cz+d_{\rm den}$ and write the printed variable as $d=-d_{\rm den}$.
The printed denominator is therefore $cz-d$, and the Fourier phase changes
with this sign choice. Thus replacing $d$ by $-d$ changes the congruence and
the phase together, not one sign in isolation.
The normalization includes the orbit conditions
\[
 (c,d)=1,\quad(c,q)=w,\quad cd\equiv uw\pmod{(w^2,q)}.
\]
With $c=w\gamma$, the last condition is
$\gamma d\equiv u\pmod{(w,q/w)}$. The cusp set has the form
\[
 u/w,\quad w\mid q,\quad u\pmod{(w,q/w)},\quad
 (u,(w,q/w))=1,
\]
with admissible lifts chosen coprime to $q$.
The sources are Iwaniec, Chapter~6 and Theorem~3.4;
DI82 (1.17), Lemmas~2.3 and~3.6;
and Kiral--Young, Proposition~3.1 and Corollary~3.2
\cite[Chapter~6; Theorem~3.4]{IwaniecSpectralMethods2002}
\cite[(1.17), p.~227; Lemma~2.3, p.~240; Lemma~3.6, p.~247]{DI82}
\cite[Proposition~3.1, Corollary~3.2]{KiralYoung2017}.
These are the conventions used in Paper~I, (2.7.T20)--(2.7.T24)
and (2.7.T41)--(2.7.T43)
\cite[(2.7.T20)--(2.7.T24), (2.7.T41)--(2.7.T43), pp.~18--19, 22--23]{PaperI}.
\end{fact}

\begin{fact}[(F13) Further spectral estimates]\label{fact:remaining-spectral-inputs}
For exceptional eigenfunctions, the first line below follows from (F3), since
$975/4096>768/4096=3/16$ and $7/32<1/2$.
Paper~I derives the interval estimate from (F5), (F6), (F7),
(P1.1) and this first line.
The last line is its supported imaginary-transform estimate:
\begin{equation}\label{eq:remaining-spectral-inputs}
 \begin{gathered}
 \lambda_j\ge3/16,\qquad 0<\vartheta_j\le1/2,\\
 \sum_{2\le q\le Q}\sum_{j\ {\rm exceptional}}^{(q)}V^{\vartheta_j}
       \left|\sum_{n\in I}\rho_j(n)\right|^2
 \ll_\epsilon(QY)^\epsilon(Q+Y+\sqrt{YV})Y,\\
 Q,Y,V\ge1,\qquad I\subset(Y,2Y]\cap\mathbb Z\quad\text{an integer interval},\\
 |\widehat\phi(i\nu)|+|\check\phi(i\nu)|
 \ll\mathfrak B(1+|\log Z|)(1+Z^{-2\nu})/(1+Z),
 \quad0<\nu\le1/4.
 \end{gathered}
\end{equation}
The last line uses the supported test class of (P1.1).
The interval argument uses DI82's proof of Theorem~7, (8.17),
with the corrected same-cusp normalization (F7)
\cite[proof of Theorem~7, (8.17)]{DI82}
\cite[(2.9), (2.7.T32), pp.~6, 20]{PaperI}.
The weaker gap $3/16$ is also DI82, Theorem~4, p.~232,
which we do not use separately.
The reconstruction of the last estimate uses (F9).
\end{fact}

\begin{fact}[(F14) Transform identities and interchanges]\label{fact:transform-interchanges}
Paper~I proves the following identities and accompanying
interchange bounds in (2.7.T26)--(2.7.T32).
Use the transforms \eqref{eq:di-transforms} without rescaling.
For $\phi\in C_c^\infty((0,\infty))$, put
\begin{equation}\label{eq:spectral-transform-applicability}
 \begin{gathered}
 \mathcal R_\phi(v)=\int_0^\infty K_{2iv}(x)\phi(x)x^{-2}\,dx,\qquad
 k_\phi(v)=\frac4{\pi^2}v\sinh(2\pi v)\mathcal R_\phi(v),\\
 \mathcal H(t,v)=\frac{\cosh\pi t}
   {\cosh\pi(t-v)\cosh\pi(t+v)},\\
 \int_{\mathbb R}k_\phi(v)K_{2iv}(x)\,dv=\phi(x)/x,\\
 \int_{\mathbb R}v\sinh(2\pi v)\mathcal H(t,v)K_{2iv}(x)\,dv
       =\cosh(\pi t)xK_{2it}(x).
 \end{gathered}
\end{equation}
By (F10), the KL hypotheses hold for both applications,
to $\phi(x)/x$ and to $xK_{2it}(x)$.
For the latter, the product-integral condition is
$|\Re(2it)|+|\Re(2iv)|<1$, with $v\in\mathbb R$.
The permitted $t$ are real or exceptional, $|\Im t|\le7/64$.
The differentiations and contour, limit, sum, and integral interchanges
in (2.7.T26)--(2.7.T32) hold in these ranges.
This includes absolute summability after the spectral estimates.
For each fixed smooth test family, the decay bound is
\[
 \mathcal R_\phi(v)\ll_{\phi,j}e^{-\pi|v|}(1+|v|)^{1-2j}
 \qquad(j\ge0).
\]
Let $\operatorname{supp}\phi\subset[Z/21,9Z]$.
The integrations by parts and differentiated estimates retain the
five-derivative seminorm $\mathfrak B$ of (P1.1), uniformly in every
auxiliary parameter.
For the following bounds $U>0$, and the oscillatory kernels are the
normalized combinations in \eqref{eq:di-transforms}; these estimates are
not assertions about an unnormalized $J_{2it}$ tail. The hyperbolic factor
in the raw Bessel representation is cancelled only after this normalization.
For truncated normalized oscillatory representations, the tails are bounded by
\[
 O\!\left(\frac{1+|t|}{x\sinh U}\right)\quad(t\in\mathbb R),\qquad
 O\!\left(\frac{(1+\nu)e^{2\nu U}}{x\sinh U}\right)
       \quad(t=i\nu,\ 0<\nu\le1/4).
\]
These tails permit the stated limiting interchanges.
The derivative and decay constants are uniform in the families
consumed by Theorem~2.8, with the losses prescribed there.
The first inversion identity follows from (F10).
The product integral, second kernel identity, and uniform
differentiation and interchange steps are proved in the same passage
\cite[(2.7.T26)--(2.7.T32), pp.~19--20]{PaperI}.
A reconstruction from the individual Bessel representations uses (F9).
\end{fact}

\begin{fact}[(P1.1) Supported real transforms]\label{fact:real-transforms}
\leavevmode\par
For $Z>0$, let $\operatorname{supp}\phi\subset[Z/21,9Z]$ and set
$\mathfrak B=\max_{0\le j\le5}\sup_{x>0}x^j|\phi^{(j)}(x)|$.
Then
\begin{equation}\label{eq:real-transform-bound}
 \begin{aligned}
 |\widehat\phi(t)|+|\check\phi(t)|
 &\ll\frac{\mathfrak B(1+|\log Z|)}{1+Z}
       \min\{1,((1+Z)/|t|)^3\},\quad t\in\mathbb R,\\
 |\widetilde\phi(l)|
 &\ll\frac{\mathfrak B(1+|\log Z|)}{1+Z}
       \min\{1,((1+Z)/l)^3\},\quad l\in\mathbb N.
 \end{aligned}
\end{equation}
The constants are absolute. The minimum at $t=0$ is $1$.
This is the real-parameter part of Paper~I, (2.7.T31)
\cite[(2.12), (2.7.T31), pp.~8, 20]{PaperI}.
The real-transform bound uses exactly the Bessel representations
in (F9). It does not cover the imaginary-parameter or interchange steps.
\end{fact}

Write $\gamma_{\!E}$ for Euler's constant. The continuous linear-sieve
functions are normalized by
\begin{equation}\label{eq:classical-delay}
 \begin{gathered}
 sF(s)=2e^{\gamma_{\!E}}\quad(0<s\le3),\qquad
 sf(s)=0\quad(0<s\le2),\\
 (sF(s))'=f(s-1)\quad(s>3),\qquad
 (sf(s))'=F(s-1)\quad(s>2).
 \end{gathered}
\end{equation}
Thus $F(1)=2e^{\gamma_{\!E}}$ and $f(s)=0$ for $0<s\le2$.
This is Paper~I, (2.23), p.~10 \cite[(2.23), p.~10]{PaperI}.

\begin{fact}[(P1.2) Structured linear sieve]\label{fact:structured-sieve}
Let $\mathcal A$ be a finite integer sequence with nonnegative
multiplicities, let $\mathcal P$ be a set of primes, and let $\mathcal X>0$.
For a multiplicative function $\omega$, set
\begin{equation}\label{eq:classical-sieve-data}
 \begin{gathered}
 P(z)=\prod_{\substack{p<z\\p\in\mathcal P}}p,\qquad
 S(\mathcal A,\mathcal P,z)=\sum_{a\in\mathcal A}1_{(a,P(z))=1},\\
 |\mathcal A_d|=\frac{\omega(d)}d\mathcal X+r(\mathcal A,d),\qquad
 V(z)=\prod_{\substack{p<z\\p\in\mathcal P}}
                       \left(1-\frac{\omega(p)}p\right).
 \end{gathered}
\end{equation}
Assume $0<\omega(p)<p$ on $\mathcal P$, $\omega(p)=0$ off $\mathcal P$,
and $\omega(p^a)\ge0$. Suppose fixed $K\ge2$ and $L_\omega>1$ satisfy,
for every $z>w\ge2$,
\begin{equation}\label{eq:classical-density}
 \begin{gathered}
 \prod_{\substack{w\le p<z\\p\in\mathcal P}}
 \left(1-\frac{\omega(p)}p\right)^{-1}
 \le\frac{\log z}{\log w}\left(1+\frac K{\log w}\right),\\
 \sum_{\substack{w\le p<z\\p\in\mathcal P}}
       \sum_{a\ge2}\frac{\omega(p^a)}{p^a}
 \le\frac{L_\omega}{\log(3w)}.
 \end{gathered}
\end{equation}
Fix $D\ge2$, $0<\epsilon_s<1/3$ and $\delta_s=\epsilon_s^9$.
Define
\begin{equation}\label{eq:classical-grid}
 \mathcal G(D)=\{D^{\epsilon_s^2(1+\delta_s)^n}:n\ge0\}.
\end{equation}
Let $\mathcal H(D)$ consist of the finite tuples
$\mathbf D=(D_1,\ldots,D_k)$ from this grid such that
$D_k\le\cdots\le D_1<D^{1/2}$. Set
\begin{equation}\label{eq:classical-tail-tuples}
 \begin{aligned}
 \mathcal D^+(D)&=\{\mathbf D\in\mathcal H(D):
 D_1\cdots D_{2l}D_{2l+1}^{3}<D\quad
             (0\le l\le(k-1)/2)\},\\
 \mathcal D^-(D)&=\{\mathbf D\in\mathcal H(D):
 D_1\cdots D_{2l-1}D_{2l}^{3}<D\quad
             (1\le l\le k/2)\}.
 \end{aligned}
\end{equation}
Each of these sets has fewer than $\exp(8\epsilon_s^{-3})$ elements.
There are coefficients $\phi_d^\nu$ and $\lambda_d^\nu(\mathbf D)$,
$\nu\in\{+,-\}$, with $|\phi_d^\nu|,|\lambda_d^\nu(\mathbf D)|\le1$, independent of
$\mathcal A,\mathcal X,\omega$, for which
\begin{equation}\label{eq:classical-structured-sieve}
 \begin{aligned}
 S(\mathcal A,\mathcal P,z)&\le \mathcal X V(z)\{F(s)+E_s\}+R_1^++R^+
                   &&(2\le z\le D),\\
 S(\mathcal A,\mathcal P,z)&\ge \mathcal X V(z)\{f(s)-E_s\}-R_1^--R^-
                   &&(2\le z\le D^{1/2}),\\
 s&=\frac{\log D}{\log z},\qquad
 E_s\ll\epsilon_s+\epsilon_s^{-8}e^{K+L_\omega}(\log D)^{-1/3}.
 \end{aligned}
\end{equation}
The implied constant in the last line is absolute. The signed remainders are
\begin{equation}\label{eq:classical-smooth-remainder}
 R_1^\nu=\sum_{\substack{d<D^{\epsilon_s}\\
                        d\mid P(D^{\epsilon_s^2})}}
                  \phi_d^\nu r(\mathcal A,d)
\end{equation}
and
\begin{equation}\label{eq:classical-tail-remainder}
 R^\nu=\sum_{\mathbf D\in\mathcal D^\nu(D)}
 \sum_{\substack{d<D^{\epsilon_s}\\d\mid P(D^{\epsilon_s^2})}}
 \lambda_d^\nu(\mathbf D)
 \sum_{\substack{D_i\le p_i<D_i^{1+\delta_s}\\p_i\mid P(z)\\
                 p_i\ne p_j\ ((i,j)\in I(\mathbf D))}}
 r(\mathcal A,dp_1\cdots p_k).
\end{equation}
Here $I(\mathbf D)$ imposes distinctness between the equal-label
positions in Iwaniec's construction. Unequal grid intervals are disjoint.
The coefficients may depend on $\mathcal P,D,z,\epsilon_s,\mathbf D$
and these distinctness conditions.
We do not claim that coefficients at different cutoffs agree. See Paper~I, Theorem~2.6,
(2.22), (2.24), (2.27)--(2.33), pp.~10--12
\cite[Theorem~2.6, (2.22), (2.24), (2.27)--(2.33), pp.~10--12]{PaperI}.
Its classical source is Iwaniec, Theorem~4
\cite[Theorem~4]{Iwaniec1980Remainder}.
\end{fact}

Below, the prime set is $\mathcal P_N=\{p:p\nmid2N\}$.
Take the nominal density $\omega(p)=1$ on $\mathcal P_N$,
$\omega(p)=0$ elsewhere and $\omega(p^a)=0$ for $a\ge2$. Then the
second sum in \eqref{eq:classical-density} is zero. Mertens' product
estimate gives one fixed absolute $K_{\mathrm M}\ge2$ with
\begin{equation}\label{eq:classical-mertens-density}
 \prod_{w\le p<z}(1-1/p)^{-1}
 \le\frac{\log z}{\log w}\left(1+\frac{K_{\mathrm M}}{\log w}\right)
 \qquad(z>w\ge2).
\end{equation}
The first product in \eqref{eq:classical-density} is no larger than this
product. Hence $K=K_{\mathrm M}$, $L_\omega=2$ are valid for
this nominal density. The factor in
\eqref{eq:classical-structured-sieve} is then $e^{K_{\mathrm M}+2}$.
See \cite[pp.~10--11]{PaperI}. Each actual sequence still needs
its own decomposition \eqref{eq:classical-sieve-data} and error estimate.

For every relative level $L_c\ge2$ we use, exactly as in Paper~I,
(2.33a)--(2.33b), p.~12 \cite[(2.33a)--(2.33b), p.~12]{PaperI},
\begin{equation}\label{eq:classical-relative-grid}
 \begin{gathered}
 0<\sigma\le10^{-10},\qquad \tau=\sigma^9,\qquad
 \epsilon_s=\sigma,\quad\delta_s=\tau,\quad D=L_c,\\
 \mathcal G(L_c)=\{L_c^{\sigma^2(1+\tau)^n}:n\ge0\},\qquad
 d<L_c^\sigma,\quad d\mid P(L_c^{\sigma^2}).
 \end{gathered}
\end{equation}
Thus the prime intervals, the smooth-divisor support and the level all
refer to the same $L_c$. First fix $\sigma>0$. Then let $N$ tend
to infinity.


\begin{fact}[(P1.3) Finite factorization]\label{fact:hb}
Let $X\ge2$, $K\in\mathbb N$, and $z\ge1$.
Assume $z^K>X$.
For every integer $1\le n\le X$,
\begin{equation}\label{eq:classical-hb}
 \Lambda(n)=\sum_{j=1}^{K}(-1)^{j-1}\binom Kj
 \sum_{\substack{m_1\cdots m_jn_1\cdots n_j=n\\m_i\le z}}
 \mu_{\rm Mob}(m_1)\cdots\mu_{\rm Mob}(m_j)\log n_1.
\end{equation}
This is Paper~I, Lemma~2.7, (2.36), p.~13
\cite[Lemma~2.7, (2.36), p.~13]{PaperI}.
\end{fact}

\begin{proof}
For $\Re s>1$, set $M_z(s)=\sum_{m\le z}\mu_{\rm Mob}(m)m^{-s}$
and $Y_z(s)=\zeta(s)M_z(s)$. The polynomial identity
$1-(1-Y_z)^K=\sum_{j=1}^K(-1)^{j-1}\binom KjY_z^j$ gives
\[
 -\frac{\zeta'}\zeta=
 \sum_{j=1}^K(-1)^{j-1}\binom Kj(-\zeta')\zeta^{j-1}M_z^j
 +\left(-\frac{\zeta'}\zeta\right)(1-Y_z)^K.
\]
Every nonzero coefficient of $Y_z-1$ has index greater than $z$.
The last term therefore has no coefficient at $n\le X<z^K$;
coefficient extraction proves \eqref{eq:classical-hb}.
\end{proof}

For $n\le2X_m$, choose $K=3$, $z=2X_m^{1/3}$.
Then $z^3=8X_m>2X_m$.
Any interval partition is imposed after \eqref{eq:classical-hb}.
The product condition and all representation multiplicities remain in
the resulting coefficients. See \cite[(2.36)--(2.37), p.~13]{PaperI}.

For the packing theorem, put $\mu_*=25/128$ and define the packing ceiling by
\[
 H_{\rm pk,max}=\frac{153}{256}.
\]
Use this constant in the packing hypotheses below. Put
\begin{equation}\label{eq:classical-packing-functions}
 \begin{gathered}
 v_{\rm pk}(x)=\min\left\{\frac{1+x}{2},\frac{32-34x}{39},
                         \frac{57}{64}-\frac32x\right\}\quad(x\ge0),\\
 \psi_{\rm pk}(x,y)=x1_{x\ge y}\quad(x\in\mathbb R,\ y\ge0),\\
 w_{\rm pk}(x_1,x_2,x_3)=\psi_{\rm pk}\left(\min\left\{
 \frac{132-68x_3}{214},\frac{153-96x_3}{256},1-2x_2\right\},\frac{1+x_1}{2}\right).
 \end{gathered}
\end{equation}
For $x_1\ge x_2\ge x_3\ge0$, define
\begin{equation}\label{eq:classical-packing-ceiling}
 \begin{aligned}
 H(x_1,x_2,x_3)=\max\{&v_{\rm pk}(x_1),v_{\rm pk}(x_2),
 v_{\rm pk}(x_1+x_2),v_{\rm pk}(x_1+x_2+x_3),\\
 &w_{\rm pk}(x_1,x_2,x_3),w_{\rm pk}(x_1,x_3,x_2),\\
 &\psi_{\rm pk}(v_{\rm pk}(x_1+x_3),x_1+2x_2+x_3),\\
 &\psi_{\rm pk}(v_{\rm pk}(x_2+x_3),2x_1+x_2+x_3)\},\\
 \widetilde H(x_1,x_2,x_3)&=
 \begin{cases}H(x_1,x_2,x_3),&x_1\le\mu_*,\\v_{\rm pk}(x_1),&x_1>\mu_*.
 \end{cases}
 \end{aligned}
\end{equation}
These are the definitions preceding Paper~I, Theorem~4.1, p.~53
\cite[p.~53]{PaperI}.

\begin{fact}[(P1.4) Three-factor packing]\label{fact:packing}
Let $0<\delta_{\rm pk}\le3\cdot10^{-9}$,
$1/2\le h\le H_{\rm pk,max}$, $d_*=h-2\delta_{\rm pk}$, and
$\delta_{\rm pk}\le g\le1/3-\delta_{\rm pk}/2$.
Let $e_1\ge\cdots\ge e_m\ge0$ be a finite labelled list.
Equal values keep distinct labels. The empty list is allowed.
Extend the list by zeros, writing $\bar e_i=0$ for $i>m$.
Suppose
\begin{equation}\label{eq:classical-packing-hypotheses}
 \sum_{i<j}e_i+2e_j\le h-2\delta_{\rm pk}\quad(1\le j\le m),\qquad
 h\le\widetilde H(\bar e_1,\bar e_2,\bar e_3).
\end{equation}
The rule of Paper~I, Section~4.1, partitions the labels as
$J_1\sqcup J_2\sqcup J_3$. With
$\Sigma_s=\sum_{i\in J_s}e_i$ and $\theta_{\rm sp}=7/32$,
\begin{equation}\label{eq:classical-packing-loads}
 \begin{aligned}
 \Sigma_1&\le g,\\
 2g+\Sigma_2+2\Sigma_3&\le1-4\delta_{\rm pk},\\
 \theta_{\rm sp}+g+\Sigma_1+5\Sigma_2+2\Sigma_3&\le2-4\delta_{\rm pk},\\
 (1+2\theta_{\rm sp})g+\Sigma_1+
 (5-2\theta_{\rm sp})\Sigma_2+
 (2+2\theta_{\rm sp})\Sigma_3
 &\le2-4(1+\theta_{\rm sp})\delta_{\rm pk}.
 \end{aligned}
\end{equation}
The partition depends only on these data.
This is Paper~I, Theorem~4.1, (4.1)--(4.3), pp.~53--54
\cite[Theorem~4.1, (4.1)--(4.3), pp.~53--54]{PaperI}.
\end{fact}


\begin{fact}[(P1.5) Continuous Hecke comparison]\label{fact:eisenstein-hecke}
Assume the three Eisenstein statements in (F12).
Let $w\mid q$, put $w'=(w,q/w)$, and let $\chi$ be a Dirichlet
character modulo $w'$.
Then Paper~I, (2.7.T41)--(2.7.T43), gives
\begin{equation}\label{eq:eisenstein-arithmetic-coefficient}
 \begin{aligned}
 \varphi_{u/w,\infty,n}(z)
  &=(w'/(wq))^z\sum_{\substack{\gamma\ge1\\(\gamma,q/w)=1}}\gamma^{-2z}
 \sum_{\substack{d\ (\gamma w)\\(d,\gamma w)=1\\\gamma d\equiv u\ (w')}}
                      e(-nd/(\gamma w)),\\
 \mathcal V_\chi(n;z)&=\sum_{u\ (w')}^*\chi(u)
                         \varphi_{u/w,\infty,n}(z),\\
 \mathcal V_\chi(kn';z)
   &=\overline{\chi(k)}\left(\sum_{d\mid k}\chi(d)^2d^{1-2z}\right)
                           \mathcal V_\chi(n';z),
       \quad(k,q)=(k,n')=1.
 \end{aligned}
\end{equation}
The first identity initially holds for $\Re z>1$.
The last extends meromorphically.
For finitely supported $v_{n'}$ in this coprime range,
\begin{equation}\label{eq:eisenstein-hecke-loss}
 \begin{aligned}
 &\sum_{\mathfrak a}\left|\sum_{n'}v_{n'}|kn'|^{it}
       \varphi_{\mathfrak a\infty,kn'}(1/2+it)\right|^2\\
 &\qquad\le d(k)^2\sum_{\mathfrak a}\left|\sum_{n'}v_{n'}|n'|^{it}
       \varphi_{\mathfrak a\infty,n'}(1/2+it)\right|^2.
 \end{aligned}
\end{equation}
It holds almost everywhere in real $t$, for either index sign
\cite[(2.7.T41)--(2.7.T43), pp.~22--23]{PaperI}.
\end{fact}

\begin{proof}
In \eqref{eq:eisenstein-arithmetic-coefficient}, write the modulus
$c_1=w\gamma=c_1'c_1''$, where $(c_1',q)=1$ and every prime of
$c_1''$ divides $q$. CRT factors the inner character sum as
$\chi(c_1')G_\chi(-n;c_1'')c_{c_1'}(n)$. Here $G_\chi$ is the
corresponding finite Gauss sum and $c_{c_1'}$ is the Ramanujan sum.
The outer character supplies another $\chi(c_1')$. The coprime factor is
\[
 L^{(q)}(2z,\chi^2)^{-1}
 \sum_{\substack{d\mid |n|\\(d,q)=1}}\chi^2(d)d^{1-2z},
\]
where $L^{(q)}$ omits Euler factors at primes dividing $q$.
For $(k,q)=(k,n')=1$, its divisor sum factors at $n=kn'$.
Changing the Gauss-sum variable supplies $\overline{\chi(k)}$.
Absolute convergence for $\Re z>1$, followed by meromorphic continuation,
proves the last identity in \eqref{eq:eisenstein-arithmetic-coefficient}.
At $z=1/2+it$, the character multiplier is
\[
 \overline{\chi(k)}\sum_{d\mid k}\chi(d)^2(k/d^2)^{it}.
\]
Its modulus is at most $d(k)$.
Orthogonality over the complete cusp set proves the inequality.
\end{proof}

\begin{fact}[(P1.6) Exceptional parameters on the full space]\label{fact:full-space-spectral-gap}
Apply (F3) to $\Gamma_0(q)$.
Every exceptional eigenfunction in the full $\Gamma_0(q)$ space then has
\[
 \vartheta_j=2\sqrt{1/4-\lambda_j}\le7/32.
\]
The bound depends only on the eigenvalue, so it covers oldforms.
These are the parameters in Paper~I, (2.7.T32), (2.7.T52), and (2.7.T56)
\cite[(2.7.T32), (2.7.T52), (2.7.T56), pp.~20, 24--25]{PaperI}.
\end{fact}

\begin{fact}[(P1.7) The spectral estimate]\label{fact:spectral-implication}
Paper~I, Theorem~2.8, proves the estimate below. Its proof uses
the large sieves (F6) and the trace formulas (F7)--(F8). It also uses the transform
bounds (P1.1), (F13)--(F14), the full-space Hecke bounds (F11), and
the continuous and exceptional comparisons (P1.5)--(P1.6).
The normalized unitary cusp transport, proved in Paper~I,
(2.7.T33)--(2.7.T40), is as follows. Choose $rw-sy=1$ and use
\begin{equation}\label{eq:unitary-cusp-transport}
 \begin{gathered}
 W_{rs}=r^{-1/2}\begin{pmatrix}r&y\\rs&rw\end{pmatrix},\qquad
 W_{rs}\Gamma_0(rs)W_{rs}^{-1}=\Gamma_0(rs),\\
 \rho_{F,1/s}(n)=e(n\bar s/r)\rho_{F\circ W_{rs},\infty}(n),\qquad
 \psi_{F,1/s}(n)=e(n\bar s/r)\psi_{F|W_{rs},\infty}(n),\\
 \varphi_{\mathfrak a,1/s,n}(z)
   =e(n\bar s/r)\varphi_{W_{rs}^{-1}\mathfrak a,\infty,n}(z),\\
 T_kF(z)=k^{-1/2}\sum_{hd=k}\sum_{j\ (d)}F((hz+j)/d),\qquad
 c_{T_kF}(n)=\sum_{h\mid(k,n)}c_F(kn/h^2).
 \end{gathered}
\end{equation}
The last line is the Maass normalization, with $(k,rs)=1$.
For holomorphic forms use the normalized coefficients of (F11).
These formulas fix the normalization of the following theorem.

Let $C,D,Y,R,S,A_T\ge1$ and $0<\epsilon_0<1/100$.
Let $a_T\in\mathbb Z$ satisfy $0<|a_T|\le A_T$.
Let $b_{n,r,s}$ be supported on
\begin{equation}\label{eq:t28-support}
 0<n\le Y,\qquad r\sim R,\qquad s\sim S,\qquad (a_T,rs)=1.
\end{equation}
Let $g$ be smooth, supported in
$(C,2C]\times(D,2D]\times\mathbb R_{>0}^3$, with
\begin{equation}\label{eq:t28-seminorms}
 |\partial^\nu g(c,d,n,r,s)|
 \le G_\nu(c^{-\nu_1}d^{-\nu_2}n^{-\nu_3}r^{-\nu_4}s^{-\nu_5})^{1-\epsilon_0}
 \quad(\nu\in\mathbb Z_{\ge0}^5).
\end{equation}
Set $\mathcal X=2+A_TCDYRS$ and $\theta_{\rm sp}=7/32$.
For every $\epsilon>0$,
\begin{equation}\label{eq:t28-conclusion}
 \left|\sum_{(r,s)=1}\sum_{0<n\le Y}b_{n,r,s}
 \sum_{\substack{c,d\ge1\\(rd,sc)=1}}
 g(c,d,n,r,s)e\left(\frac{a_T n\overline{rd}}{sc}\right)\right|
 \ll_{\epsilon,\epsilon_0,\{G_\nu\}}
       \mathcal X^{\epsilon+18\epsilon_0}J_*(a_T),
\end{equation}
where the inverse is modulo $sc$, and
\begin{equation}\label{eq:t28-main-bound}
 \begin{aligned}
 J_*(a_T)^2={}&\sum_{\substack{n''\mid a_T^\infty\\n''\le2Y}}
 A_T^{\theta_{\rm sp}}(n'')^{\theta_{\rm sp}}\\
 &\quad\times\left[CS\left(\frac Y{n''}+RS\right)(C+DR)+|a_T|YRS\right]
                 \|\widetilde b(n'')\|_2^2\\
 &+\left[CS(CD\sqrt R)^{2\theta_{\rm sp}}(Y+RS)^{1-\theta_{\rm sp}}\right.\\
 &\qquad\left.\times(C+DR)^{1-2\theta_{\rm sp}}+D^2YR\right]\|b\|_2^2,\\
 \|b\|_2^2={}&\sum_{n,r,s}|b_{n,r,s}|^2,\qquad
 \|\widetilde b(n'')\|_2^2=
       \sum_{\substack{n',r,s\\(n',a_T)=1}}|b_{n'n'',r,s}|^2.
 \end{aligned}
\end{equation}
Here $n''\mid a_T^\infty$ means that every prime factor of $n''$ divides $a_T$.
It includes $n''=1$; when $a_T=\pm1$, only $n''=1$ occurs.
The bound is uniform in $(a_T,A_T,C,D,Y,R,S,b)$.
These are Paper~I, Theorem~2.8, (2.38)--(2.43)
\cite[Theorem~2.8, (2.38)--(2.43), pp.~14--15]{PaperI}.
\end{fact}

Section~\ref{sec:l3} checks the masked coefficients, intervals and support
of the fourth switched row.

